\section{Search in a Rotated Sorted Array I}
\label{sec:bs-rotated-1}

\problemstatement{An array that was originally sorted in increasing order
has been \textbf{rotated} at some unknown pivot -- for example,
$[0,1,2,4,5,6,7]$ rotated at index $3$ becomes $[4,5,6,7,0,1,2]$. All
elements are distinct. Given such a rotated array $\textit{arr}$ and a
target $X$, find the index of $X$, or return $-1$ if it is not present.
We must still achieve $O(\log n)$ time.}

\begin{patternbox}
\emph{``Rotated sorted array''} is a clear signal that plain binary search
(Section~\ref{sec:bs-find-x}) cannot be applied directly, because the array
as a whole is no longer globally sorted -- comparing
$\textit{arr}[\textit{mid}]$ to $X$ no longer tells us which half to
discard. But the keyword \emph{``sorted, then rotated''} carries a crucial
piece of remaining structure: although the \emph{whole} array is not
sorted, \textbf{at least one of the two halves split at any midpoint always
is}. Whenever a problem mentions rotation, immediately look for this
``one half is always sorted'' structure, since it is what restores binary
search's monotonic-discard guarantee.
\end{patternbox}

\subsection*{Understanding the problem carefully}

Pick any midpoint $\textit{mid}$ in a rotated sorted array. The two halves
$[\textit{lo}, \textit{mid}]$ and $[\textit{mid}, \textit{hi}]$ cannot both
contain the rotation pivot (the single place where a larger element is
immediately followed by a smaller one) -- the pivot lies in at most one of
the two halves, so the other half is guaranteed to be a genuinely sorted,
unrotated run of increasing values. We can detect which half is sorted with
one comparison: if $\textit{arr}[\textit{lo}] \le
\textit{arr}[\textit{mid}]$, the left half is sorted; otherwise, the right
half must be sorted.

\begin{ideabox}[Greedy Decision at Each Step]
At each step, first determine which half is sorted using
$\textit{arr}[\textit{lo}] \le \textit{arr}[\textit{mid}]$. Then, check
whether $X$ could possibly lie within that sorted half's value range (using
ordinary comparisons, since that half really is sorted). If $X$'s value
falls within the sorted half's range, discard the other half and recurse
into the sorted half; otherwise, $X$ (if present at all) must be in the
\emph{other} half, so discard the sorted half and recurse into the
unsorted-looking half instead -- which, on the next iteration, will itself
be split into one sorted and one possibly-unsorted piece, maintaining the
invariant.
\end{ideabox}

\subsection*{Why this discard rule is correct}

The key invariant maintained throughout is that the current search window
$[\textit{lo}, \textit{hi}]$ is itself a rotated sorted array (a
sub-array of the original, which inherits the ``rotated sorted'' structure
since removing elements from one end of a rotated sorted run leaves a
smaller rotated sorted run). At every step we identify a genuinely sorted
half by comparing its endpoints directly -- a valid comparison because that
half truly has no rotation inside it. If $X$ lies within that sorted half's
numeric range (i.e.\ between its smallest and largest values), it can only
be found there, since the other half's values lie entirely outside this
range (otherwise the sortedness of the identified half, combined with the
single-rotation-point structure, would be violated). So discarding the
half that provably cannot contain $X$ is always safe, exactly mirroring the
discard argument from ordinary binary search, just with one extra
case-split to first identify which half is safe to reason about directly.

\subsection*{Algorithm}

\begin{enumerate}
  \item Initialize $\textit{lo} \gets 0$, $\textit{hi} \gets n-1$.
  \item While $\textit{lo} \le \textit{hi}$:
  \begin{enumerate}
    \item $\textit{mid} \gets \textit{lo}+(\textit{hi}-\textit{lo})/2$.
    \item If $\textit{arr}[\textit{mid}] = X$: return \textit{mid}.
    \item If $\textit{arr}[\textit{lo}] \le \textit{arr}[\textit{mid}]$
          (left half is sorted):
    \begin{itemize}
      \item If $\textit{arr}[\textit{lo}] \le X < \textit{arr}[\textit{mid}]$:
            search left, $\textit{hi} \gets \textit{mid}-1$.
      \item Else: search right, $\textit{lo} \gets \textit{mid}+1$.
    \end{itemize}
    \item Otherwise (right half is sorted):
    \begin{itemize}
      \item If $\textit{arr}[\textit{mid}] < X \le \textit{arr}[\textit{hi}]$:
            search right, $\textit{lo} \gets \textit{mid}+1$.
      \item Else: search left, $\textit{hi} \gets \textit{mid}-1$.
    \end{itemize}
  \end{enumerate}
  \item If the loop ends without returning, return $-1$.
\end{enumerate}

\begin{algorithm}[H]
\caption{Search in Rotated Sorted Array (Distinct Elements)}
\begin{algorithmic}[1]
\Procedure{SearchRotated}{$\textit{arr}, X$}
  \State $\textit{lo} \gets 0$, $\textit{hi} \gets n-1$
  \While{$\textit{lo} \le \textit{hi}$}
    \State $\textit{mid} \gets \textit{lo}+(\textit{hi}-\textit{lo})/2$
    \If{$\textit{arr}[\textit{mid}] = X$} \Return \textit{mid} \EndIf
    \If{$\textit{arr}[\textit{lo}] \le \textit{arr}[\textit{mid}]$}
      \If{$\textit{arr}[\textit{lo}] \le X$ \textbf{and} $X < \textit{arr}[\textit{mid}]$}
        \State $\textit{hi} \gets \textit{mid} - 1$
      \Else
        \State $\textit{lo} \gets \textit{mid} + 1$
      \EndIf
    \Else
      \If{$\textit{arr}[\textit{mid}] < X$ \textbf{and} $X \le \textit{arr}[\textit{hi}]$}
        \State $\textit{lo} \gets \textit{mid} + 1$
      \Else
        \State $\textit{hi} \gets \textit{mid} - 1$
      \EndIf
    \EndIf
  \EndWhile
  \State \Return $-1$
\EndProcedure
\end{algorithmic}
\end{algorithm}

\subsection*{Dry run}

\dryrun{Take $\textit{arr} = [4,5,6,7,0,1,2]$, $X = 0$.}

\begin{itemize}
  \item $\textit{lo}=0,\textit{hi}=6$: $\textit{mid}=3$,
        $\textit{arr}[3]=7 \ne 0$. Is $\textit{arr}[0]=4 \le
        \textit{arr}[3]=7$? Yes, left half sorted. Is
        $4 \le 0 < 7$? No ($0 < 4$). So search right:
        $\textit{lo}=4$.
  \item $\textit{lo}=4,\textit{hi}=6$: $\textit{mid}=5$,
        $\textit{arr}[5]=1 \ne 0$. Is $\textit{arr}[4]=0 \le
        \textit{arr}[5]=1$? Yes, left half sorted. Is
        $0 \le 0 < 1$? Yes. Search left: $\textit{hi}=4$.
  \item $\textit{lo}=4,\textit{hi}=4$: $\textit{mid}=4$,
        $\textit{arr}[4]=0 = 0$. Return $4$.
\end{itemize}

The answer is index $4$.

\subsection*{C++ implementation}

\begin{lstlisting}
#include <bits/stdc++.h>
using namespace std;

int searchRotated(vector<int>& arr, int X) {
    int n = arr.size();
    int lo = 0, hi = n - 1;

    while (lo <= hi) {
        int mid = lo + (hi - lo) / 2;
        if (arr[mid] == X) return mid;

        if (arr[lo] <= arr[mid]) { // left half sorted
            if (arr[lo] <= X && X < arr[mid]) {
                hi = mid - 1;
            } else {
                lo = mid + 1;
            }
        } else { // right half sorted
            if (arr[mid] < X && X <= arr[hi]) {
                lo = mid + 1;
            } else {
                hi = mid - 1;
            }
        }
    }
    return -1;
}

int main() {
    vector<int> arr = {4, 5, 6, 7, 0, 1, 2};
    cout << "Index of 0: " << searchRotated(arr, 0) << endl;
    return 0;
}
\end{lstlisting}

\begin{complexitybox}
\textbf{Time Complexity.} Each iteration discards at least one full half
of the current window, exactly as in ordinary binary search, giving
\[
  O(\log n).
\]
\textbf{Space Complexity.} $O(1)$ auxiliary space.
\end{complexitybox}

\begin{takeawaybox}
When an array's global sortedness is broken in one specific, structured way
(here, a single rotation point), look for a local property that still
holds at every step of binary search -- here, ``at least one of the two
halves around any midpoint is genuinely sorted.'' Identifying and exploiting
that local invariant is what lets binary search survive structural
violations of global sortedness.
\end{takeawaybox}
